\documentclass[10pt,letterpaper]{article}
\usepackage{fontspec}
\setmainfont{Latin Modern Roman}[SmallCapsFont={Latin Modern Roman Caps}]
\usepackage[margin=0.8in,top=0.55in,bottom=0.55in]{geometry}
\usepackage{titlesec}
\usepackage{enumitem}
\usepackage[hidelinks]{hyperref}
\pagestyle{empty}
\setlength{\parindent}{0pt}
\setlength{\parskip}{5pt}
\urlstyle{same}
\titleformat{\section}{\large\bfseries}{}{0pt}{}[\vspace{-6pt}\rule{\textwidth}{0.5pt}]
\titlespacing*{\section}{0pt}{10pt}{6pt}
\setlist[itemize]{leftmargin=1.4em,itemsep=3pt,topsep=2pt,parsep=0pt,label=\textbullet}

\begin{document}

\begin{center}
{\LARGE\bfseries\textsc{Jose Alberto Tupayachi Silva}}\\[4pt]
{\small 661-365-5289 $\cdot$ \href{mailto:jtupayac@vols.utk.edu}{jtupayac@vols.utk.edu}; \href{mailto:tupayachisja@ornl.gov}{tupayachisja@ornl.gov}\\
\href{https://www.linkedin.com/in/jose-tupayachi-582339147}{linkedin.com/in/jose-tupayachi-582339147} $\cdot$ \href{https://github.com/jtupayachi}{github.com/jtupayachi} $\cdot$ \href{https://jtupayachi.github.io/}{jtupayachi.github.io}}
\end{center}

\section{Teaching Philosophy}

I teach industrial and systems engineering as a discipline of making decisions with data under real constraints. My own path, from process analytics in industry to building AI systems at a national laboratory to graduate research in digital twins and optimization, has convinced me that students learn these ideas best when they use them on problems that resemble the ones they will meet after graduation. My goal is for students to leave a course able to frame a problem, choose and apply the right method, and explain the result to someone who did not build the model.

\textbf{Learning by doing, anchored in real examples.} I introduce a concept through a concrete case before formalizing it. In my guest lecture on digital twins, I used intermodal transportation as the running use case so that abstract ideas such as simulation, optimization, and real-time data integration were tied to decisions students could picture. In my Git workshop, I taught version control through the day-to-day situations researchers actually face (collaborating on a shared project, recovering from a mistake, tracking what changed) rather than as a list of commands.

\textbf{Bringing research and industry into the classroom.} My work spans freight logistics, mobile health applications, geospatial analytics, and machine learning pipelines, and I have worked on teams that deliver software and models in Agile sprints. I draw on these experiences to show students how methods behave outside textbook settings: data are messy, requirements change, and results must be reproducible and communicated clearly. I want students to see research tools such as version control, scripting, and reproducible workflows as everyday professional habits, not extras.

\textbf{Assessment that supports learning.} Through my teaching assistant roles I have helped prepare exams and homework and assess student work. I design assessments to match what I ask students to do in class: problems that require reasoning and application, not only recall, paired with clear expectations and timely feedback. I view feedback as part of instruction, so students can see how to improve before the next assignment.

\textbf{An accessible, collaborative environment.} Students arrive with different backgrounds, from programming-heavy to purely analytical. I try to lower the barrier to entry by explaining tools in plain terms, offering practical starting points, and encouraging teamwork on projects. Having studied and worked in more than one country and across disciplines, I value classrooms where different perspectives help the whole group solve a problem better.

\textbf{Teaching the current toolkit responsibly.} AI, optimization, and simulation tools are changing how engineers work. I want students to use them with understanding: to know what a model assumes, how to validate it, and where it can fail. I aim to prepare students to adapt as tools change by grounding them in fundamentals, statistical thinking, and good engineering judgment.

\section{Teaching Experience}

I have supported instructional delivery, exam and homework preparation, and student assessment through my teaching assistant roles at the University of Tennessee, Knoxville (UTK). In addition, I have experience delivering technical presentations and workshops:

\begin{itemize}
\item \textbf{Guest Lecture, University of Central Florida.} Delivered a guest lecture on digital twins in collaboration with Professor Soheil Sabri, presenting concepts and applications of digital twins with intermodal transportation use cases.
\item \textbf{Git Workshop, University of Tennessee, Knoxville.} Delivered a workshop for a research methods class, helping students learn version control and research collaboration tools through day-to-day examples.
\end{itemize}

\end{document}
